\documentclass[../main.tex]{subfiles}
\graphicspath{{\subfix{../assets/}}}

\begin{document}
\subsection{Engineered Ceramics}
\subsubsection{Oxide Ceramics}
\paragraph{Alumina}
Synthetically produced, alumina (Al$_2$O$_3$) is the most common engineering ceramic and typically a first choice when considering options.
Raw alumina powder may be calcinated to form ceramics with well-known properties: electrical insulation, chemical resistance, and extreme hardness. \parencite{aluminabritannica, generalceramics}
Some images of alumina are included below.

\begin{figure}[h!]
	\begin{center}
		\begin{subfigure}[t]{0.42\textwidth}
			\includegraphics[width=\textwidth]{alumina1.jpg}
			\caption{Alumina microstructure \parencite{aluminabritannica}}
		\end{subfigure}%
		~~
		\begin{subfigure}[t]{0.32\textwidth}
			\includegraphics[width=\textwidth]{alumina2.jpg}
			\caption{Some alumina parts \parencite{photodump}}
		\end{subfigure}
		\vskip 0.2cm
		\begin{subfigure}[h]{0.5\textwidth}
			\includegraphics[width=\textwidth]{alumina3.jpg}
			\caption{Sodium vapour lamp with alumina enclosure \parencite{aluminabritannica}}
		\end{subfigure}
		\caption{Alumina examples}
	\end{center}
\end{figure}

Some applications of alumina include:
\begin{itemize}
	\item Spark plug insulation
	\item Bone implants
	\item Abrasive tools
	\item Industrial refractory linings
\end{itemize}
\parencite{aluminabritannica}

\paragraph{Zirconia}
Similarly to alumina, zirconia (ZrO$_2$) is another chemically resistant ceramic with high hardness, but differs in toughness -- while alumina is brittle, zirconia instead exhibits a remarkable resistance to crack propagation and fracturing.
Although zirconia is considered an engineering ceramic, it's likely better known in its \textit{cubic} form, resembling diamonds.
While the aforementioned mechanical properties and extremely high melting point (2700$^\circ$C) are indeed attractive, unstabilised zirconia rapidly degrades at higher temperatures, nullifying any value it brings.
Some examples are shown below:

\begin{figure}[h!]
	\begin{center}
		\begin{subfigure}[t]{0.45\textwidth}
			\includegraphics[width=\textwidth]{zirconia1.jpg}
			\caption{Zirconia plates}
		\end{subfigure}%
		~~
		\begin{subfigure}[t]{0.45\textwidth}
			\includegraphics[width=\textwidth]{zirconia2.jpg}
			\caption{Zirconia rings}
		\end{subfigure}
		\caption{Some zirconia examples \parencite{photodump}}
	\end{center}
\end{figure}

Similarly, the major applications of zirconia include:
\begin{itemize}
	\item Dental implants
	\item Jewellery
	\item Temperature sensors
	\item Thermal barriers
\end{itemize}
\parencite{zirconia}
\newpage

\subsubsection{Carbide Ceramics}
\paragraph{Silicon carbide}
A common trend among engineering ceramics is their extraordinary hardness; silicon carbide (SiC) is second only to zirconia out of the common engineering ceramics. \parencite{generalceramics}
Silicon carbide is also considered a semiconductor, particularly in IC (integrated circuit) boards and LED (light-emitting-diode) substrates. \parencite{sicbritannica}
Below is an example of a SiC crystal.

\begin{figure}[h!]
	\begin{center}
		\includegraphics[width=0.75\textwidth]{siliconcarbide.jpg}
	\end{center}
	\caption{A silicon carbide crystal (ruler for scale) \parencite{sicbritannica}}
\end{figure}

As previously mentioned, applications include:
\begin{itemize}
	\item Abrasive tools
	\item Semiconductors
	\item Refractory linings
	\item Military armour
\end{itemize}
\parencite{sicbritannica, generalceramics}
\newpage

\paragraph{Tungsten carbide}
The industry standard for cutting tools and industrial machines, tungsten carbide (WC) is another sintered ceramic finding widespread use in a diverse range of industries.
While it presents as a grey powder in its pure form, it may be sintered or cemented with carrier components to create tools. \parencite{wc}
Tungsten carbide is often bonded with cobalt (Co) during sintering to raise its density and toughness, but at a proportional cost to hardness. \parencite{generalceramics}
Below is a property table of WC alongside WC-tipped drill bits.

\begin{figure}[h!]
	\begin{center}
		\includegraphics[width=0.75\textwidth]{WC.png}
	\end{center}
	\caption{Tungsten carbide fact sheet \& WC-tipped drill bits \parencite{wc}}
\end{figure}

While tool-making is the most common and often first thought of application of WC, some other applications include:
\begin{itemize}
	\item Mining drills
	\item Military weapons
	\item Aerospace components
\end{itemize}
\newpage

\subsubsection{Nitride Ceramics}
\paragraph{Aluminium nitride}
The exceptional thermal shock resistance and conductivity of aluminium nitride (AlN) combined with its low thermal expansion ratio grant it a unique suitability for use in aerospace and electrical engineering.
Manufacturing is similar to other engineering ceramics; AlN powder is sintered with a carrier to form parts.

\begin{figure}[h!]
	\begin{center}
		\begin{subfigure}[t]{0.45\textwidth}
			\includegraphics[width=\textwidth]{AlN1.jpg}
			\caption{AlN plates}
		\end{subfigure}%
		~~
		\begin{subfigure}[t]{0.45\textwidth}
			\includegraphics[width=\textwidth]{AlN2.jpg}
			\caption{AlN disks and plates}
		\end{subfigure}
	\end{center}
	\caption{Aluminium nitride ceramics \parencite{photodump}}
\end{figure}

Often deemed an indispensable element of modern electronics, some applications include:
\begin{itemize}
	\item Heat sinks and baseplates
	\item Inter-component insulation
	\item LED substrates
	\item Piezoelectric devices
\end{itemize}
\parencite{alunitride}
\newpage

\subsubsection{Forming Methods of Ceramics}
The manufacturing processes of ceramics may be generalised into a few main stages:
\begin{itemize}
	\item Material preparation
	\item Forming (moulding)
	\item Sintering
	\item Post-processing
\end{itemize}

Within this forming stage, there are generally 3 options: slip casting, extrusion, and pressing. \parencite{forming1, forming2}
Diagrams of all 3 are included for clarity.

\begin{figure}[h!]
	\begin{center}
		\begin{subfigure}[t]{0.4\textwidth}
			\includegraphics[width=\textwidth]{slip1.jpg}
			\caption{Slip casting, part 1 \parencite{slipcast}}
		\end{subfigure}%
		~~
		\begin{subfigure}[t]{0.4\textwidth}
			\includegraphics[width=\textwidth]{slip2.jpg}
			\caption{Slip casting, part 2 \parencite{slipcast}}
		\end{subfigure}

		\begin{subfigure}[t]{0.55\textwidth}
			\includegraphics[width=\textwidth]{extrusionceramics.jpg}
			\caption{Extrusion diagram \parencite{extrudingceramic}}
		\end{subfigure}%
		~~
		\begin{subfigure}[t]{0.28\textwidth}
			\includegraphics[width=\textwidth]{pressing.jpg}
			\caption{Ceramic pressing diagram \parencite{pressingceramic}}
		\end{subfigure}
	\end{center}
	\caption{A collection of diagrams}
\end{figure}
\newpage
\paragraph{Slip casting}
In essence, this is as simple as filling a mould with liquid clay.
A high level process would be as follows:
\begin{enumerate}
	\item Prepare the mould \& slip
	\item Pour the slip into the mould \& leave the desired length of time necessary
	\item Drain the remaining slip \& wait until leather-hard
	\item Demould the part \& process as necessary
\end{enumerate}
(Simplified from \textcite{slipcast, slipcast2})

With reference to \textbf{Figure 23:a, 23:b}, steps 1-3 similarly show this process. Step 4 combines \textbf{Figure 23:b} into a singular step.

\paragraph{Extruding}
Similar to metallic extrusion as discussed in \textbf{Section 4.2.1}, ceramics (and by extension, clay) can be formed in the same way.
While differing slightly in the use of a screw rather than a hydraulic ram to drive the material, the process remains largely the same as compared to the aforementioned metallic process.
However, a standard process is still provided:
\begin{enumerate}
	\item Material preparation
	\item Loading the extruding machine
	\item Extrude the material through dies to form the desired shape
	\item Once dry, trim to desired length \& await further processing
\end{enumerate}

The extruding machine is depicted in \textbf{Figure 23:c}.

\paragraph{Pressing}
Perhaps the simplest out of the 3 methods, pressing is as the name suggests: press the material into the mould as hard as necessary to form the desired shape.
Although there are 3 types of pressing (uniaxial, isostatic, wet), the principle is identical for all:
\begin{enumerate}
	\item Load the mould with material
	\item Apply pressure until desired density is reached
	\item Wait until touch-dry, then remove and post-process
\end{enumerate}

Isostatic pressing is illustrated in \textbf{Figure 23:d}.
\newpage

\subsection{Analysis of AI}
The ChatGPT response without edits:

\fbox{%
	\parbox{\dimexpr\linewidth-2\fboxsep-2\fboxrule}{%
		\vskip 0.5cm
		\itshape
		Method for Slip Casting a Ceramic Hand Basin
		\begin{enumerate}
			\item Prepare the Mould: Clean and assemble a two-part plaster mould, ensuring all surfaces are smooth and dry. Apply a release agent, such as soapy water, to ease demoulding.
			\item Pour the Slip: Fill the mould with liquid clay (slip) to the brim. Allow it to sit for 15–30 minutes to form a consistent wall thickness as the plaster absorbs moisture.
			\item Drain Excess Slip: Pour out the remaining slip, leaving a hollow shell inside the mould. Let it set until leather-hard.
			\item Demould the Basin: Carefully remove the mould parts without damaging the cast. Smooth imperfections.
			\item Dry and Fire: Let the basin air-dry completely before bisque firing and applying glazes for final firing.
		\end{enumerate}
		\rmfamily
		\vskip 0.5cm
		\parencite{chatgpt}
	}%
}
\vskip 1cm

Perhaps the most glaring omission in the response is the preparation of the slip.
The lack of this detail in itself is sufficient to show the unlikely success of this method, as it is contingent on the quality of the slip and the mould, neither of which have given technical detail.
In particular, professional guides describe a \textit{deflocculation} process that's entirely missed by ChatGPT.

As a result, consequences of product quality may include:
\begin{itemize}
	\item Failure of the part
	\item Failure of the slip to dry or slip dries too quickly
	\item Poor aesthetics
\end{itemize}

\textcite{slipcast2} mentions the ratio of slip components as the following, noting that this slip neither formed a film nor settled overnight, indicating correct deflocculation:
\begin{itemize}
	\item 10000g porcelain powder
	\item 8.8 litres water
	\item 44g \textit{Darvan} deflocculant
\end{itemize}

Furthermore, \textcite{slipcast2} also stresses the importance of proper slip viscosity, noting that 100ml of slip should drain from a 4.25 mm hole in approximately 40 seconds.
\newpage
\begin{figure}[h!]
	\begin{center}
		\begin{subfigure}[t]{0.45\textwidth}
			\includegraphics[width=\textwidth]{slipbucket.jpg}
			\caption{Correctly deflocculated slip}
		\end{subfigure}%
		~~
		\begin{subfigure}[t]{0.45\textwidth}
			\includegraphics[width=\textwidth]{slipcup.jpg}
			\caption{Viscosity measuring cup}
		\end{subfigure}
	\end{center}
	\caption{Images from a guide by \textcite{slipcast2}}
\end{figure}

\begin{figure}[h!]
	\begin{center}
		\begin{subfigure}[t]{0.55\textwidth}
			\includegraphics[width=\textwidth]{baddeflocculation.jpg}
			\caption{Incorrect deflocculation}
		\end{subfigure}%
		~~
		\begin{subfigure}[t]{0.35\textwidth}
			\includegraphics[width=\textwidth]{dontmakeplates.jpg}
			\caption{Poor results with flat moulds}
		\end{subfigure}
	\end{center}
	\caption{Examples of mistakes from a guide by \textcite{slipcast2}}
\end{figure}

Hence, the success of this product is entirely contingent on the quality of the slip and the mould; ChatGPT's failure to provide this technical information results in a determination of poor quality.
\newpage
Hence, the revised response:

\fbox{%
	\parbox{\dimexpr\linewidth-2\fboxsep-2\fboxrule}{%
		\vskip 0.5cm
		\itshape
		Method for Slip Casting a Ceramic Hand Basin
		\begin{enumerate}
			\item \textbf{Prepare slip with a 1.8:1 ratio of porcelain powder to water}
			\item \textbf{Add deflocculant until the viscosity is such that 100ml of slip should drain from a 4.25 mm hole in approximately 40 seconds (approximately 1g of deflocculant per 100ml of water)}
			\item Prepare the Mould: Clean and assemble a two-part plaster mould, ensuring all surfaces are smooth and dry. Apply a release agent, such as soapy water, to ease demoulding.
      \item \textbf{Test slip drying time and determine the desired length of time to leave the slip in the mould}
      \item Pour the Slip: Fill the mould with liquid clay (slip) to the brim. Allow it to sit for \textbf{the length of time previously determined} to form a consistent wall thickness as the plaster absorbs moisture.
			\item Drain Excess Slip: Pour out the remaining slip, leaving a hollow shell inside the mould. Let it set until leather-hard; \textbf{ensure this is not left for longer than necessary as cracking may occur as seen in Figure25:d}.
			\item Demould the Basin: Carefully remove the mould parts without damaging the cast. Smooth imperfections.
			\item Dry and Fire: Let the basin air-dry completely before bisque firing and applying glazes for final firing.
		\end{enumerate}
		\rmfamily
		\vskip 0.5cm
	}%
}

As previously presented, \textbf{Figure 23:a and 23:b} are excellent graphical examples of the procedure.
\newpage

\sloppy
\printbibliography[heading=subbibnumbered,keyword={sect3},title={References}]
\fussy
\end{document}
